\chapter{Erste Schritte, Systemgrenzen \& Quick Start} \label{ch:quickstart}

\section{Grundverständnis \& Systemgrenzen} \label{sec:systemgrenzen}

\begin{tcolorbox}[colback=warnbg,colframe=warnred,arc=2mm,title=\textbf{ACHTUNG: Theoretische Simulation vs. Reale Fertigung}]
Die Berechnungen dieses Tools basieren auf analytischen Modellen unter Annahme idealer Materialeigenschaften. 
Die realen Prüfstandsergebnisse weichen in Abhängigkeit von der \textbf{Fertigungsqualität und äußeren Einflüssen} ab.
\end{tcolorbox}

\begin{itemize}
    \item \textbf{Einfluss der Fertigungsqualität}: Faktoren wie der effektive Anpressdruck im Heißpressverfahren, Temperaturzyklen, Faserauslenkung, Porengehalt und Harzaufnahme der offenen Wabenkammern beeinflussen die Steifigkeit und Festigkeit maßgeblich.
    \item \textbf{Erwartete Abweichung}: Eine Abweichung von \textbf{ca. 10\,\%} zwischen Simulation und Prüfstandsergebnis ist bei sorgfältiger, reproduzierbarer Prepreg-Fertigung normal und als Auslegungskorridor zu berücksichtigen.
    \item \textbf{Nichtlinearität der Energieaufnahme}: Während der elastische Anstieg Gradient sehr präzise vorhergesagt werden kann, unterliegt die Energieberechnung bis zur Normdurchbiegung von $12{,}7\,\text{mm}$ größeren Streuungen. 
        Grund hierfür ist der hochgradig nichtlinearität im gesammten Bereich vor und nach Überschreiten der Maximalkraft. Zur fundierten Plausibilisierung der absorbierten Energie wurden diskrete Versagensmodi implementiert (weiterführende Details siehe Technische Dokumentation).   
\end{itemize}
\clearpage










\section{Kurzbeschreibung \& Systemzweck} \label{sec:systemzweck}

Das Tool dient der schnellen, unkomplizierten und regelkonformen Vorauslegung von Faserverbund-Sandwichplatten (\hyperref[abk:CFK]{CFK}-Deckschichten mit Aluminium-Wabenkern oder Hartschaum) für das Formula Student Monocoque.

Es bildet den genormten \textbf{3-Punkt-Biegeversuch} (nach \hyperref[abk:ASTM]{ASTM} C393 / D7250) rechnerisch ab und ermittelt die für das \hyperref[abk:SES]{\textbf{Structural Equivalency Spreadsheet (SES)}} geforderten Strukturkennwerte:
\begin{enumerate}
    \item Äquivalente Biegesteifigkeit (\hyperref[abk:EI]{$EI$}) des Verbunds
    \item Elastisches Biegemodul und Steifigkeitsgradient
    \item Maximale Bruchkraft ($F_{max}$) unter Berücksichtigung Versagensmodi
    \item Energieabsorption ($E_{12.7}$) bis $12{,}7\,\text{mm}$ Stempelweg
    \item und weitere informative Größen
\end{enumerate}










\section{Repository klonen mit Sourcetree} \label{sec:sourcetree}

Um die Versionsverwaltung über Git durch eine übersichtliche und intuitiv bedienbare grafische Benutzeroberfläche (\hyperref[abk:GUI]{GUI}) zu vereinfachen, wird die Verwendung von Atlassian Sourcetree empfohlen:

\begin{enumerate}
    \item \textbf{Insatlation}: https://www.atlassian.com/software/sourcetree
    \item \textbf{Optional GitHub Account}: https://github.com 
    \begin{itemize}
        \item \textbf{Favorisieren}: Favorisiere das Repository um es leichter wieder zu finden.
    \end{itemize}
    \item \textbf{Repository-URL}: https://github.com/Basso-RS/SES-Panel-Berechnung.git 
    \item \textbf{Sourcetree öffnen}: Klicke oben auf die Registerkarte \textbf{„Klonen / Neu“} (\textit{Clone}).
    \item \textbf{Pfade angeben}:
    \begin{itemize}
        \item \textbf{Quellpfad / URL}: Füge die kopierte Git-URL ein.
        \item \textbf{Zielfeld}: Wähle dein lokales Arbeitsverzeichnis.
    \end{itemize}
    \item \textbf{Klonen ausführen}: Klicke auf \textbf{„Klonen“}. Sourcetree lädt alle Dateien, Skripte und Materialdatenbanken herunter.
    \item \textbf{Branch prüfen}: Stelle sicher, dass du auf dem Haupt-Branch (\texttt{main} bzw. \texttt{master}) arbeitest.
\end{enumerate}
\clearpage









\section{Quick Start: In 4 Schritten zur ersten Berechnung} \label{sec:quickstart}

Das gesamte Pfad- und Umgebungsmanagement ist automatisiert! 
\par In vier Schritten führst du deine erste Simulation durch:





\subsection*{Schritt 1: Projekt in MATLAB initialisieren}

Öffne MATLAB und navigiere in das geklonte Projektverzeichnis. \par
Die Projektdatei \texttt{SES\_Panelberechnung.prj} im Verzeichnisbaum  von MATLAB (Links) per Doppelklick oder Rechtsklick ausführen:

\textit{(Hintergrund: Durch das Öffnen des MATLAB-Projekts werden alle erforderlichen Unterordner (\texttt{00\_Scripts}, \texttt{01\_Functions}, \texttt{02\_Data}) automatisch in den MATLAB-Suchpfad aufgenommen.)}





\subsection*{Schritt 2: Panel \& Lagenaufbau prüfen}
Öffne die Datei \texttt{01\_Functions/Scenariodata.m}. \par
Hier ist der zu berechnende Prüfkörper definiert:
\begin{lstlisting}[style=vscodeLightMatlab]
Scenario.panel.b     = 275;   % mm Panel-Breite
Scenario.panel.l     = 400;   % mm Stützweite
Scenario.panel.w     = 12.7;  % mm Max. Durchbiegung
\end{lstlisting}

Layup Informationen für eine Berechnung eines Panels:

\begin{lstlisting}[style=vscodeLightMatlab]
Scenario.name   = '26_16';
Scenario.core_t = 18;        % 18 mm Wabenkern
Scenario.layup  = '45t7/0m/0m/0m/21m/45t7/cp/18core/cp/45t7/-21m/0m/0m/0m/45t7';
\end{lstlisting}
Passe bei Bedarf den Namen, die Kerndicke (\texttt{core\_t}) oder den Layup-String an und speichere die Datei.

Zu beachten ist hier die angegebene Abkürzungskonvention der Material- und Lagen-Tokens:
\begin{itemize}
    \item \textbf{Faserlagen}: Zusammensetzung aus \texttt{[Winkel][Kürzel]}, z.\,B. \texttt{45t7} ($+45^\circ$ T700-Gewebe) oder \texttt{-21m} ($-21^\circ$ \hyperref[abk:UD]{M40J-UD}).
    \item \textbf{Wabenkern}: Deklaration über \texttt{[Dicke]core}, z.\,B. \texttt{18core} für $18\,\text{mm}$ Aluwabe.
    \item \textbf{Klebefilm}: \hyperref[abk:CP]{\texttt{cp}} deklariert den beidseitig am Kern verwendeten Harzfilm (CP006).
\end{itemize}

\begin{table}[htbp]
\centering
\small
\begin{tabular}{llcp{7.5cm}}
\toprule
\textbf{Kürzel} & \textbf{Materialbezeichnung} & \textbf{Typ} & \textbf{Funktion / Eigenschaften} \\
\midrule
\texttt{t7}   & T700 Fabric         & Gewebe & Zähe Sichtlage ($45^\circ$) \& Kernanbindung \\
\texttt{m}    & M40J UD             & UD     & Hauptlastaufnahme ($0^\circ, \pm 21^\circ$, Hochmodul) \\
\texttt{f}    & M40J Fabric         & Gewebe & Hochsteife Decklage (Hochmodul-Gewebe) \\
\texttt{t8}   & T800S UD            & UD     & Hochfeste Zwischenlagen (IM-Faser) \\
\texttt{z}    & ZZ300T              & Gewebe & Zylon/PBO für maximale Energieaufnahme \\
\texttt{a}    & SIGAPREG A W173     & Gewebe & Aramid/Kevlar für Durchschlagschutz \\
\texttt{cp}   & Coupling Film CP006 & Harz   & Klebefilm zwischen Kern und Deckhaut \\
\texttt{core} & PAMG 5056 Aluwabe   & Kern   & Wabenkern (z.\,B. \texttt{18core}) \\
\bottomrule
\end{tabular}
\caption{Materialkürzel und Basisfunktionen im Lagenaufbau.}
\label{tab:quickstart_materials}
\end{table}


\subsection*{Schritt 3: Skript ausführen}

Führe das Berechnungsskript \texttt{Calculate\_Single\_Panel}, 
aus dem Ordner \texttt{00\_Scripts} über RUN im Editor-Tab aus. 

\textit{(Alternativ: Datei \texttt{00\_Scripts/Calculate\_Single\_Panel.m} im MATLAB-Editor öffnen und \texttt{F5} drücken).}





\subsection*{Schritt 4: Ergebnisse einsehen}
Nach Abschluss der Berechnung (Rechenzeit typischerweise $< 0{,}1\,\text{Sekunden}$):
\begin{enumerate}
    \item \textbf{Terminal-Report}: Im MATLAB Command Window erscheint eine detaillierte Zusammenfassung aller Steifigkeits-, Festigkeits- und Kennwerte.
    \item \textbf{Automatischer PDF-Prüfbericht}: Das Skript erzeugt automatisch einen DIN-A4-Prüfbericht und legt diesen mit Zeitstempel im persönlichen \textbf{\texttt{Downloads}}-Verzeichnis ab:
    \begin{quote}
        \texttt{SES\_Report\_Single\_Panel\_26\_20\_<Zeitstempel>.pdf}
    \end{quote}
\end{enumerate}
